\section{Overview}\label{sec:introduction:overview}
Quantum computing hardware has matured to the point where hybrid quantum-classical programs can be compiled and executed on real devices.
Gate-based quantum processors now support circuits of enough depth and qubit count to execute non-trivial quantum algorithms on current hardware.
Yet the programming models through which researchers and engineers interact with this hardware have not kept pace.
Every existing quantum programming system, from low-level circuit assemblers to data structure frameworks operating entirely in superposition~\cite{yuan2025foundationalabstractions}, requires the programmer to reason in quantum-mechanical terms.
Even where data structure operations are presented through a classical-facing API, quantum-mechanical obligations persist at the programming surface.
Reversibility must be managed explicitly, uncomputation must be scheduled by the programmer and quantum-computational primitives remain beyond the abstraction boundary.
The physical constraints of the underlying hardware must inform programming decisions at every level of the program.
No existing system provides a surface model that is fully \textit{agnostic} of the underlying quantum substrate~\cite{DiMatteo2024AbstractionHierarchy,nunez_corrales2025betterabstractmachines}.
The abstraction boundary that separates quantum-mechanical implementation concerns from programmer-facing vocabulary has not been fully crossed.

This is not a usability concern only.
Decoherence, the physical process by which quantum states lose coherence due to environmental interaction, imposes hard temporal limits on circuit execution.
Quantum hardware is characterised by coherence timescales $T_1$ (how long a non $\ket{0}$ state can be maintained) and $T_2$ (how long a state in superposition remains stable), which bound the period over which a qubit's state remains reliable.
Circuits whose depth exceeds the hardware's coherence window produce results that are statistically indistinguishable from noise.
Unnecessarily deep circuits, produced by compilers that lack the semantic information to optimise at a higher level of abstraction, are therefore more likely to produce physically incorrect results on current hardware.
The programming model is a direct determinant of circuit quality: a language that forces explicit circuit construction cannot, in general, supply the compiler with the information needed to produce shorter, more coherent circuits automatically.
Just as a classical programmer can separate their intent for a program from the process and memory management that executes it, a quantum programmer should not require knowledge of the underlying quantum substrate.
No existing language has delivered this separation in full.

The work presented in this thesis demonstrates that this separation is formally achievable.